% included from 5conclusion.tex
\section{Discussion and Limitations}
\label{sec:limits}

\paragraph{What the evidence supports} (i)~Conditioning a graph-attention forecaster on the timetable reduces load-forecast error substantially over the six-interval horizon, in particular in the announced surge windows (Section~\ref{sec:exp-pred}), although at the one-step horizon simple forecasters are more accurate; the effect is robust over seeds and is not obtained by the graph or Transformer layers without the schedule. (ii)~Per-AP SAC agents started from the incumbent configuration remain close to it, whereas the DDPG-based learned baselines trained with the same protocol drift away and lose reward, mostly through channel switching (Section~\ref{sec:exp-main}). (iii)~Ignoring MAC contention and channel-switching cost changes the evaluation materially (Section~\ref{sec:exp-robust}), which is a caution for simulator-only studies of Wi-Fi resource management, including our first version.

\paragraph{What it does not support} The learned controllers do not outperform the static incumbent plan in the simulator, although a perfect-information optimiser shows that headroom exists. Neither feeding the ST-GCAT embedding into the policy state (Contribution~2) nor the bidirectional joint gradient (Contribution~3) improves control or forecasting in our experiments; the joint gradient with $\eta=0.1$ made both slightly worse. These two components should be regarded as unproven. The reward weights are design choices, not derived quantities (Section~\ref{sec:exp-weights}).

\paragraph{Simulator} The simulator is quasi-static (one evaluation per 5-minute interval) and is not calibrated to measurements; no measured traces are used. Hidden terminals, adjacent-channel interference, frame-level dynamics, rate adaptation and client roaming (users always associate to the nearest AP) are not modelled; contention uses the saturation model of Bianchi with a fixed share of backlogged stations. The timetable is synthetic and exact. The channel-switching outage ($3$~s) and the EDCA airtime mapping are engineering assumptions that we did not validate on hardware; the fourth EDCA category and vendor-specific airtime-quota interfaces are not covered. Results are for one venue layout with 30 APs.

\paragraph{Training and statistics} Each run uses 1000 episodes (200 of exploration around the incumbent configuration); longer training, other hyperparameters or other exploration schemes might change the control results. The number of seeds is small (three for the main comparison), so we report effect sizes and bootstrap intervals; no result of the paper depends on a significance threshold. The design choices made while developing the method (reward scale, initial entropy temperature, update-to-data ratio, initialisation at the incumbent configuration, embedding scale $\kappa$) were selected on validation episodes of one seed, mostly with variants of STAG-FedSAC; the baselines were not tuned separately, which favours the proposed method over the baselines. All learned methods start from the incumbent configuration, which is an inductive bias; it favours stability and may not be available when no sensible configuration exists.

\paragraph{Privacy and federation} Federation is used for data minimisation, not for a formal guarantee: neither differential privacy nor secure aggregation is applied, and model parameters and distilled outputs can leak information. The predictor is centralised at the venue controller and consumes AP-level aggregates. The public reference set is generated from an independent proxy venue, which assumes that such a simulator is available to all parties; its refresh traffic is included in the measured communication.

\paragraph{Future work} Packet-level simulation (e.g.\ ns-3~\cite{ns3ai}) with hidden terminals and rate adaptation; calibration and validation on measured venue traces such as~\cite{dartmouth2018}; recomputing embeddings for stored transitions (or learning with fresh embeddings) to make the joint gradient usable; longer training and safe-improvement schemes that let the learner depart from the incumbent only where it has evidence of gains; delayed or uncertain timetables; and differential privacy or secure aggregation for the federation.

\section{Conclusion}
\label{sec:conclusion}

We studied schedule-driven WiFi resource management in transit venues with a schedule-conditioned graph-attention forecaster, per-AP Soft Actor-Critic agents and a two-level federation, in a simulator that, unlike the first version of this paper, models MAC contention, channel-switching cost and a realisable QoS mapping. The clearest result is on forecasting: conditioning on the timetable lowers the RMSE by \redFullPersistRmse\% relative to persistence and \redFullATwoRmse\% relative to the same predictor without the schedule, and by \redFullATwoSurgeRmse\% in surge windows. On control, the learned federated agents are stable and far better than learned DDPG baselines, but they do not outperform the static incumbent configuration in our simulator, and we could not show a benefit of feeding the embedding to the policy or of the joint gradient. Federation improves on isolated agents (reward $\mFullReward$ against $\mIsoReward$), and exchanging soft policy outputs between zones needs $\commRatioFlat\times$ less wide-area traffic than flat averaging and $\commRatioAFive\times$ less than weight-based hierarchical federation at a similar reward; the reward differences among the federated variants are within seed noise. The conclusions are limited to simulation; the open questions are listed above. All numbers in this paper are generated by the code and run records in the public repository (\url{https://github.com/vinhqdang/STAG_FedSAC}); no trained checkpoints are distributed.
